%-----------------------------------------------------------------------------------------------------------------------------------------------%
%	The MIT License (MIT)
%
%	Copyright (c) 2021 Jitin Nair
%
%	Permission is hereby granted, free of charge, to any person obtaining a copy
%	of this software and associated documentation files (the "Software"), to deal
%	in the Software without restriction, including without limitation the rights
%	to use, copy, modify, merge, publish, distribute, sublicense, and/or sell
%	copies of the Software, and to permit persons to whom the Software is
%	furnished to do so, subject to the following conditions:
%	
%	THE SOFTWARE IS PROVIDED "AS IS", WITHOUT WARRANTY OF ANY KIND, EXPRESS OR
%	IMPLIED, INCLUDING BUT NOT LIMITED TO THE WARRANTIES OF MERCHANTABILITY,
%	FITNESS FOR A PARTICULAR PURPOSE AND NONINFRINGEMENT. IN NO EVENT SHALL THE
%	AUTHORS OR COPYRIGHT HOLDERS BE LIABLE FOR ANY CLAIM, DAMAGES OR OTHER
%	LIABILITY, WHETHER IN AN ACTION OF CONTRACT, TORT OR OTHERWISE, ARISING FROM,
%	OUT OF OR IN CONNECTION WITH THE SOFTWARE OR THE USE OR OTHER DEALINGS IN
%	THE SOFTWARE.
%	
%
%-----------------------------------------------------------------------------------------------------------------------------------------------%

%----------------------------------------------------------------------------------------
%	DOCUMENT DEFINITION
%----------------------------------------------------------------------------------------

% article class because we want to fully customize the page and not use a cv template
\documentclass[a4paper,11pt]{article}

%----------------------------------------------------------------------------------------
%	FONT
%----------------------------------------------------------------------------------------

% % fontspec allows you to use TTF/OTF fonts directly
% \usepackage{fontspec}
% \defaultfontfeatures{Ligatures=TeX}

% % modified for ShareLaTeX use
% \setmainfont[
% SmallCapsFont = Fontin-SmallCaps.otf,
% BoldFont = Fontin-Bold.otf,
% ItalicFont = Fontin-Italic.otf
% ]
% {Fontin.otf}

%----------------------------------------------------------------------------------------
%	PACKAGES
%----------------------------------------------------------------------------------------
\usepackage{url}
\usepackage{parskip} 	

%other packages for formatting
\RequirePackage{color}
\RequirePackage{graphicx}
\usepackage[usenames,dvipsnames]{xcolor}
\usepackage[scale=0.94]{geometry}

%tabularx environment
\usepackage{tabularx}

%for lists within experience section
\usepackage{enumitem}

% centered version of 'X' col. type
\newcolumntype{C}{>{\centering\arraybackslash}X} 

%to prevent spillover of tabular into next pages
\usepackage{supertabular}
\usepackage{tabularx}
\newlength{\fullcollw}
\setlength{\fullcollw}{0.47\textwidth}

%custom \section
\usepackage{titlesec}				
\usepackage{multicol}
\usepackage{multirow}
%custom bullet points - 
\usepackage{enumitem}

\setlist[itemize]{nosep,after=\strut,leftmargin=1em,itemsep=3pt,label=\textendash}

%CV Sections inspired by: 
%http://stefano.italians.nl/archives/26
\titleformat{\section}{\Large\scshape\raggedright}{}{0em}{}[\titlerule]
\titlespacing{\section}{0pt}{10pt}{10pt}

%for publications
\usepackage[style=authoryear,sorting=ynt, maxbibnames=2]{biblatex}

%Setup hyperref package, and colours for links
\usepackage[unicode, draft=false]{hyperref}
\definecolor{linkcolour}{rgb}{0,0.2,0.6}
\hypersetup{colorlinks,breaklinks,urlcolor=linkcolour,linkcolor=linkcolour}
\addbibresource{citations.bib}
\setlength\bibitemsep{1em}

%for social icons
\usepackage{fontawesome5}

% Define a command for the link icon and optionally redefine \href
\newcommand{\linkicon}{{\footnotesize\faExternalLink*}}
% Uncomment the following line to automatically add the icon to all \href links
% \let\oldhref\href
% \renewcommand{\href}[2]{\oldhref{#1}{#2\linkicon}}

%debug page outer frames
%\usepackage{showframe}


% job listing environments
\newenvironment{jobshort}[2]
    {
    \begin{tabularx}{\linewidth}{@{}l X r@{}}
    \textbf{#1} & \hfill &  #2 \\[3.75pt]
    \end{tabularx}
    }
    {
    }

\newenvironment{joblong}[2]
    {
    \begin{tabularx}{\linewidth}{@{}l X r@{}}
    \textbf{#1} & \hfill &  #2 \\[3.75pt]
    \end{tabularx}
    \begin{minipage}[t]{\linewidth}
    \begin{itemize}[nosep,after=\strut, leftmargin=1em, itemsep=3pt,label=--]
    }
    {
    \end{itemize}
    \end{minipage}    
    }



%----------------------------------------------------------------------------------------
%   BEGIN DOCUMENT
%----------------------------------------------------------------------------------------
\begin{document}

\pagestyle{empty} 

%----------------------------------------------------------------------------------------
%   TITLE & CONTACT
%----------------------------------------------------------------------------------------
\begin{tabularx}{\linewidth}{@{} C @{}}
\Huge{Mend-Amar Badral} \\[7.5pt]
\href{https://github.com/MendeBadra}{\raisebox{-0.05\height}\faGithub\ MendeBadra} \ $|$ \ 
\href{https://www.linkedin.com/in/mend-amar-badral-64217a28a/}{\raisebox{-0.05\height}\faLinkedin\ Mend-Amar Badral} \ $|$ \ 
\href{https://mendebadra.github.io/}{\raisebox{-0.05\height}\faGlobe \ mendebadra.github.io} \\
\href{mailto:mendamar.badral@gmail.com}{\raisebox{-0.05\height}\faEnvelope \ mendamar.badral@gmail.com} \ $|$ \ 
\href{tel:+36 707385744}{\raisebox{-0.05\height}\faMobile \ +36 70 7385 744} \ $|$ \
\faMapMarker \ Budapest, Hungary 
\end{tabularx}

%----------------------------------------------------------------------------------------
%   PERSONAL STATEMENT - Tailored for DTA Startup Focus
%----------------------------------------------------------------------------------------
% \vspace{1em}
% \noindent
% A Computer Science student with a strong foundation in Applied Mathematics, driven by a philosophy of continuous, incremental improvement. I'm on a mission to bridge the gap between theoretical results and functional software, with a focus on Computer Vision and 3D technologies.
% I'm a diligent and hard-working student with a passion for applying computational thinking to the most interesting of problems. I aspire to contribute to the advancements in scientific discovery, public understanding of mathematics and open-source technologies. I'm an avid Julia enthusiast and Python programmer.

% Diligent and hard-working student with a passion of computational thinking and software. Always aspiring to contribute to the creation of value, I believe in the power of baby-steps and importance of teamwork.


% I'm currently a MSc Computer Science student with a BSc in Applied Mathematics and a focus on Machine Learning. I'm passionate about applying statistical techniques and algorithmic problem-solving to high-impact quantitative finance challenges.

% I'm a diligent and hard-working student with a passion for facing new and exciting challenges. I believe in the power of baby steps and collaborative spirit of people.

% \noindent

% A mathematician-turned-engineer with a passion for building software from first principles. I specialize in bridging the gap between abstract logic and scalable implementation. Driven by intellectual stamina, a collaborative spirit, and a commitment to solving complex technical challenges through rigorous habit and teamwork.


% Computer Science M.Sc.\ student with a background in Applied Mathematics and a strong interest in reliable software systems. Experienced in building automation scripts, working in Linux environments, and integrating simulation and data processing pipelines. I value disciplined work habits, steady technical improvement, and collaborative approach to problem solving.

% Personal statement
Driven by a problem solving mindset and fueled by curiosity, I aspire to develop genuine solutions for the real-world, impactful problems. 
%----------------------------------------------------------------------------------------
%   SKILLS
%----------------------------------------------------------------------------------------
\vspace{-0.5em}
\section{Skills}
\noindent
\begin{tabularx}{\linewidth}{@{}l X@{}}
\textbf{Languages:} & English (C1 \href{https://drive.google.com/file/d/1IxH-0NrfMyf-ZqEpQFMrmtcXiDdzc00k/view?usp=sharing}{IELTS 8.0 \linkicon}) \\

\textbf{Programming:} & Python, SQL, Bash, Julia, C/C++, R\\

\textbf{Tools:} & Linux, Git, Docker, pandas, PyTorch, sklearn, opencv, open3d, Quarto, \LaTeX \\ % TODO: Add Kubernetes here. 


% \textbf{Data Engineering:} & Large-scale data processing, relational databases (MySQL/PostgreSQL), file formats (GeoTIFF, HDF5, Parquet) \\

\end{tabularx}

%----------------------------------------------------------------------------------------

% EDUCATION

%----------------------------------------------------------------------------------------
\vspace{-0.5em}
\section{Education}

\begin{tabularx}{\linewidth}{@{}l X r@{}}

\textbf{Budapest University of Technology and Economics (BME)} & & Sep, 2025 -- June, 2027\\

M.Sc.\ in Computer Science Engineering & & Budapest, Hungary \\

\textit{Supervisor:} \href{https://scholar.google.com/citations?user=Sns05HgAAAAJ\&hl=en}{M\'arton Vaitkus \linkicon} \\

\textit{Specialization:} Data Science \& Artificial Intelligence \\


& & \\



\textbf{National University of Mongolia (NUM)} & & Sep, 2021 -- June, 2025 \\

B.Sc.\ in Applied Mathematics & & Ulaanbaatar, Mongolia \\
GPA: 3.2/4.0 \textbf{(rank 3/34)}\\
\textit{Supervisor:} \href{https://scholar.google.com/citations?user=oy6Pcz4AAAAJ\&hl=en}{Galtbayar Artbazar \linkicon} \\

\textit{Thesis:} \href{https://heyzine.com/flip-book/ba3096dda3.html}{\textbf{Evaluating land degradation using image processing methods \linkicon}} \\
% \textit{\href{https://docs.google.com/presentation/d/1BvcVU_O69c-oM05ckK0No2ZJ_5Wgpw5ufpk8lP6x6ao/edit?usp=sharing}{Presentation}} 
% \href{https://docs.google.com/presentation/d/1BvcVU_O69c-oM05ckK0No2ZJ_5Wgpw5ufpk8lP6x6ao/edit?usp=sharing}{Presented} the results in \\ 
% \textit{The First NUM-Yonsei university Student Workshop on Mathematics and Computing}


\end{tabularx}

%----------------------------------------------------------------------------------------

% EXPERIENCE SECTIONS

%----------------------------------------------------------------------------------------
\vspace{-0.5em}
\section{Experience}
\begin{jobshort}{Center of Mathematics Application, NUM}{Oct, 2023 -- June, 2025}
\textit{Undergraduate Research Intern} \hfill 
\begin{itemize}[nosep,after=\strut,leftmargin=1em,itemsep=3pt,label=--]
    % Image processing
    \item \textbf{Data Processing \& Python:} Developed custom \href{https://github.com/MendeBadra/BelceeriinDoroitol}{\textbf{Python pipelines \linkicon}} multi-threading to process large-scale datasets (1,800+ multispectral images). Optimized data ingestion and processing workflows, reducing execution time by 80\% (5x factor).
    % GPU
    \item \textbf{Data Infrastructure:} Orchestrated photogrammetry pipelines using \textbf{OpenDroneMap} and \textbf{Docker} with \textbf{GPU acceleration} on Linux. Managed complex dependency environments and GPU debugging to ensure reliable high-throughout processing of spatial data.
    % SQL
    \item \textbf{Data Management:} Utilized \textbf{SQL} and relational database principles to organize research metadata and spatial results, facilitating efficient querying and data integrity across multiple field studies.
    
    
    % \item \textbf{Documentation \& Teamwork:} Authored documentation, and operational procedures for research workflows, improving onboarding and collaboration within a student group.


    

    % \item \textbf{System administration:} Maintained remote access to laboratory servers and computers using Tailscale and SSH, managed and distributed SSH keys, configured automated backups.


   

    % \item \textbf{Technology Evaluation:} Integrated Dockerized OpenDroneMap workflows as an open-source alternative to proprietary drone-processing software, working extensively in Linux environments with version control systems.


% \item  
% Managed and installed Docker containerized applications that is essential for the research.

% \item \textbf{Research:} Thesis research reached a conclusion that the spatial clustering of weed can be captured using Moran's I coefficient performed on drone images. 
% Enabled fast research iteration cycles by analyzing and performing comparisons on the results for discussion with the supervisor for feedback. 
% Presented the results in \textit{The First NUM-Yonsei university Student Workshop on Mathematics and Computing}. 

% \item \textbf{Seminar:} Participated in \href{https://site.num.edu.mn/dep/am/news/39}{Study Group Mathematics with Industry} where we worked on the problem of \href{https://github.com/mendeBadra/RailRoadSim/}{simulating railroad cargo delivery process} of Ulaanbaatar Railway company for optimizing costs, which revealed a possibility to achieve up to 23\% reduction in costs.

\end{itemize} 

\vspace{-1.5em}
\end{jobshort}

% Authored a lab document for other students to onboard.
% Enabled fast research iteration cycles by analyzing and performing comparisons on the results for discussion with the supervisor for feedback. 
% Presented the results in \textit{The First NUM-Yonsei university Student Workshop on Mathematics and Computing}. 
%  ensuring reliable batch processing and correctness of the output

% \item \textbf{Seminar:} Participated in \href{https://site.num.edu.mn/dep/am/news/39}{Study Group Mathematics with Industry} where we worked on the problem of \href{https://github.com/mendeBadra/RailRoadSim/}{simulating railroad cargo delivery process} of Ulaanbaatar Railway company for optimizing costs, which revealed a possibility to achieve up to 23\% reduction in costs.

%----------------------------------------------------------------------------------------
% PROJECTS
%----------------------------------------------------------------------------------------
\vspace{-0.5em}
\section{Projects}
\vspace{-0.5em}
\href{https://github.com/MendeBadra/aisim2nerfstudio}
{\textbf{aiSim Data Integration for 3D Gaussian Splatting \linkicon}}

\begin{itemize}[nosep,after=\strut,leftmargin=1em,itemsep=3pt,label=--]

  \item Developed a sensor fusion pipeline to parse and convert over sensor measurements from aiMotive's \href{https://aimotive.com/aisim}{aiSim \linkicon} ADAS simulation platform into the \href{https://docs.nerf.studio/}{nerfstudio \linkicon} 3D scene representation format. Worked on the camera coordinate system transformation from sensor frame to the world frame. 

  % \item Worked with structured JSON parsing, large-scale sensor data processing, and reproducible Linux-based development workflows.

\end{itemize}

\vspace{-1.5em}
\href{https://mendebadra.github.io/posts/bike-sharing-dataset/bike_sharing_dataset_eda.html}
{\textbf{Exploratory Data Analysis -- Bike Sharing Dataset
\linkicon}}
\begin{itemize}[nosep,after=\strut,leftmargin=1em,itemsep=3pt,label=--]
  \item Performed exploratory data analysis on bike sharing dataset. Published the results personal blog website using Quarto.
\end{itemize}


\vspace{-0.5em}
\href{https://github.com/MendeBadra/DeepLearningBishopFigures}
{\textbf{Deep Learning Foundations and Concepts Book Examples Reproduction \linkicon}}

\begin{itemize}[nosep,after=\strut,leftmargin=1em,itemsep=3pt,label=--]
  \item Reproduced polynomial regression experiments from Bishop \& Bishop's \textit{Deep Learning Foundations and Concepts} using Julia notebooks to using interactive sliders.
  % \item Generated synthetic datasets and fit models of varying polynomial degree 
  % to analyze how model complexity affects tests and training error.

  % \item Built interactive visualizations in \textbf{Julia (Pluto.jl)} to illustrate 
  % regression curves, training error, and model behavior across degrees.
\end{itemize}
% PROJECT HUNGARIAN MONGOLIAN EDA
% \noindent
% \href{https://github.com/MendeBadra/trustworthy_ai_course}{\textbf{Comparative EDA: Housing Affordability in Hungary vs Mongolia}} 

% \begin{itemize}[nosep,after=\strut,leftmargin=1em,itemsep=3pt,label=--]
%   \item Performed end-to-end exploratory data analysis (EDA) on public datasets from the 
%     \href{https://www.ksh.hu/?lang=en}{Hungarian Central Statistical Office (KSH)} and 
%     \href{https://www.1212.mn/en}{National Statistics Office of Mongolia} to determine:
%     \textit{“How many m\textsuperscript{2} of a new apartment could a monthly average net wage buy, by region and year?”}
% \end{itemize}

\noindent
% \href{https://github.com/MendeBadra/DeepLearningBishopFigures}
% {\textbf{Deep Learning Foundations and Concepts Book Examples Reproduction}}
% \begin{itemize}[nosep,after=\strut,leftmargin=1em,itemsep=3pt,label=--]
% %   \item Self-initiated project to reproduce key examples from C.~Bishop and H.~Bishop’s \textit{Deep Learning Foundations and Concepts} to gain theoretical and fundamental understanding of \textbf{polynomial regression} and \textbf{probability theory}.
%   \item Reproduced examples and figures from C.~Bishop's \textit{Deep Learning Foundations and Concepts} book.
% % \item Published reproducible Jupyter and \href{https://plutojl.org/}{Pluto.jl} notebooks on GitHub, 
% % emphasizing clear explanation, interactivity, and visualization.
% \end{itemize}

%   Streamlined conversion from simulation output to Gaussian Splatting training using Python.


% }% \\

% \end{tabularx}

\vspace{-2.5em}

% \noindent
\section{Awards}
% \begin{itemize}[nosep,after=\strut,leftmargin=1em,itemsep=3pt,label=--]
 \textbf{Stipendium Hungaricum Scholarship} \hfill Sep, 2025 
\end{document}
